\section*{Bab \ref{ch:b2:cell-differentiation} --- Diferensiasi Sel: Sel Otot Rangka}

\begin{solution}{exo:b2:cell-differentiation:1}
\omterm{def:b2:cell-differentiation:differentiation}{Diferensiasi}: pengungkapan yang mantap atas gen dan fungsi satu tipe
sel. \omterm{def:b2:cell-differentiation:differentiation}{Determinasi}: pemancangan pada sebuah nasib sebelum ia
diungkapkan. \omterm{def:b2:cell-differentiation:differentiation}{Potensi}: rentang nasib yang masih terbuka. Zigotnya:
totipoten; sel \omterm{def:b2:mammal-reproduction:implantation}{massa sel dalamnya}: pluripoten; sel satelitnya:
unipoten (otot saja) --- tetap saja sebuah \omterm{def:b2:cell-differentiation:differentiation}{sel punca}; \omterm{def:b2:cell-differentiation:myogenesis}{serabut ototnya}:
terdiferensiasi, pascamitosis, tanpa \omterm{def:b2:cell-differentiation:differentiation}{potensi}.
\end{solution}

\begin{solution}{exo:b2:cell-differentiation:2}
Sel dermomiotom (Pax3) $\to$ \omterm{def:b2:cell-differentiation:myogenesis}{mioblas} terimbas (\omterm{def:b2:cell-differentiation:myogenesis}{MyoD}, Myf5), yang
membelah selama FGF-nya bertahan $\to$ keluar dari daurnya, miogenin
$\to$ pembarisan dan peleburan menjadi sebuah \omterm{def:b2:cell-differentiation:myogenesis}{miotuba} $\to$ gen ototnya
menyala dan \omterm{def:b2:cell-differentiation:fibre}{miofibril} dirakit (MRF4) $\to$ serabut matang berinti tepi,
beserta sel satelit yang diam (Pax7).
\end{solution}

\begin{solution}{exo:b2:cell-differentiation:3}
Bahwa nukleus sebuah sel yang terdiferensiasi mempertahankan tiap gen
yang dibutuhkan untuk membuat seekor hewan utuh, dan bahwa sitoplasma
telurnya dapat menyetel ulang programnya: \omterm{def:b2:cell-differentiation:differentiation}{diferensiasi} dapat balik dan
bukan sebuah kehilangan DNA. Ia tidak menunjukkan bahwa tiap nukleus
yang terdiferensiasi dapat disetel ulang (sebagian besar pemindahannya
gagal), maupun bagaimana penyetelan ulangnya bekerja, maupun bahwa sel
yang terdiferensiasi itu sendiri dapat berganti nasib di tempat.
\end{solution}

\begin{solution}{exo:b2:cell-differentiation:4}
Cakram Z yang membatasi \omterm{def:b2:cell-differentiation:fibre}{sarkomernya}; filamen tipis aktin yang tertambat
padanya; filamen tebal miosin yang berpusat pada garis M-nya; pita A
(filamen tebal) dan pita I (tipis saja). Pada pemendekannya, filamen
tipisnya meluncur ke arah garis M-nya dan cakram Z-nya saling
mendekat: pita I dan zona H-nya menyempit, pita A --- yaitu panjang
filamen tebalnya --- tidak berubah, dan tidak ada filamen yang memendek.
\end{solution}

\begin{solution}{exo:b2:cell-differentiation:5}
$0.2/2.5\times 10^{-6} = \num{80000}$ \omterm{def:b2:cell-differentiation:fibre}{sarkomer}; $0.2/25\times 10^{-6}
= \num{8000}$ inti; satu \omterm{def:b2:cell-differentiation:myogenesis}{mioblas} tiap intinya, jadi kira-kira
\num{8000} \omterm{def:b2:cell-differentiation:myogenesis}{mioblas} yang melebur (sel satelit menambahkan lagi kemudian).
\end{solution}

\begin{solution}{exo:b2:cell-differentiation:6}
FGF mengimbas Id, yang mengikat protein pasangan yang dibutuhkan \omterm{def:b2:cell-differentiation:myogenesis}{MyoD}
untuk membentuk sebuah dimer pengikat DNA; \omterm{def:b2:cell-differentiation:myogenesis}{MyoD} hadir tetapi
menganggur, dan selnya terus membelah. Ketika FGF-nya ditarik, Id-nya
meluruh, \omterm{def:b2:cell-differentiation:myogenesis}{MyoD} berdimer, miogeninnya dinyalakan, dan selnya keluar dari
daurnya lalu melebur. \omterm{def:b2:cell-differentiation:myogenesis}{Mioblas} yang membuat Id secara tetap tidak pernah
melebur: mereka membelah tanpa batas dan tidak membentuk otot.
\end{solution}

\begin{solution}{exo:b2:cell-differentiation:7}
Syaratnya $0.4m = m^{2}/(1 + m^{2}) + 0.02$. Pada $m = 2.1$: kiri
$0.84$, kanan $4.41/5.41 + 0.02 = 0.835$: sebuah keadaan tunak. Keadaan
rendahnya: bagi $m$ yang kecil, $0.4m \approx m^{2} + 0.02$, sehingga
$m \approx
0.055$ (kiri $0.022$, kanan $0.023$). Dinaikkan ke 1, yang
berada di atas ambangnya di dekat $0.41$, selnya memanjat ke $2.1$ lalu
tinggal di situ: ia terpancang.
\end{solution}

\begin{solution}{exo:b2:cell-differentiation:8}
Dengan $k = 0.6$: pada $m = 1$, penguraiannya $0.6$ melampaui
sintesisnya $0.52$; pada $m = 0.5$, $0.30 > 0.22$; pada $m = 2$,
$1.2 > 0.82$: garisnya terletak di atas kurvanya di mana pun di luar
keadaan rendahnya, yang menjadi satu-satunya keadaan tunaknya. Sebuah
sel yang \omterm{def:b2:cell-differentiation:myogenesis}{MyoD-nya} terurai terlalu cepat tidak dapat mempertahankan
keadaan nyalanya: ia kembali ke keadaan tak terdiferensiasinya apa pun
pengimbasannya --- tanpa otot.
\end{solution}

\begin{solution}{exo:b2:cell-differentiation:9}
(a) Tanpa \omterm{def:b2:cell-differentiation:myogenesis}{mioblas}, tanpa otot rangka: kedua faktornya berlebihan bagi
\omterm{def:b2:cell-differentiation:differentiation}{determinasinya}, sehingga keduanya harus hilang. (b) \omterm{def:b2:cell-differentiation:myogenesis}{Mioblasnya}
terbentuk dan berbaris tetapi tidak pernah melebur maupun membuat
protein otot: miogenin dibutuhkan bagi \omterm{def:b2:cell-differentiation:differentiation}{diferensiasinya}. (c) Nyaris
normal, karena Myf5 menggantikan \omterm{def:b2:cell-differentiation:myogenesis}{MyoD} dalam \omterm{def:b2:cell-differentiation:differentiation}{determinasinya}.
\end{solution}

\begin{solution}{exo:b2:cell-differentiation:10}
Di dalam sebuah fibroblas, gen ototnya terletak di dalam kromatin yang
dapat dibuka \omterm{def:b2:cell-differentiation:myogenesis}{MyoD}; di dalam sebuah sel hati, gennya mungkin terjejali
di dalam heterokromatin, atau \omterm{prop:b2:cell-differentiation:transcriptional}{pengatur induk} hatinya sendiri mungkin
menekannya, sehingga selnya tidak berkompeten. Ujinya: rawatlah sel
hatinya dengan sebuah obat yang melonggarkan kromatinnya (sebuah
penghambat histon deasetilase) sebelum menambahkan \omterm{def:b2:cell-differentiation:myogenesis}{MyoD}, atau
tambahkan \omterm{def:b2:cell-differentiation:myogenesis}{MyoD} bersama sebuah faktor yang menggeser program hatinya,
lalu nilailah pengungkapan miosinnya.
\end{solution}

\begin{solution}{exo:b2:cell-differentiation:11}
Picuan berfrekuensi rendah yang bertahan menaikkan kalsium di dalam
selnya selama waktu yang lama; lintasan yang diaktifkan kalsium
menyalakan gen miosin lambatnya, biogenesis mitokondrianya,
mioglobinnya, dan pertumbuhan kapilernya, lalu mematikan isoform
cepatnya. Program induk serabutnya (keadaan \omterm{def:b2:cell-differentiation:myogenesis}{MyoD-nya}, arsitektur
\omterm{def:b2:cell-differentiation:fibre}{sarkomernya}) tidak tersentuh: latihan mengubah gen otot mana yang
diungkapkan, bukan jati diri selnya, dan tidak ada isyarat dari sebuah
saraf yang dapat membuka kembali program rawan yang telah ditutup pada
\omterm{def:b2:cell-differentiation:differentiation}{determinasinya}.
\end{solution}

\begin{solution}{exo:b2:cell-differentiation:12}
Keadaan nyala sebuah faktor yang mengaktifkan dirinya sendiri adalah
sebuah titik mantap sebuah gelung umpan balik, yang dipertahankan
sintesis yang sinambung, bukan oleh perubahan apa pun pada gennya;
tanda kromatinnya menambahkan lapisan kedua yang menjaga gen yang
dibungkam tetap tertutup. Maka penyetelan ulangnya menuntut sirkuitnya
didorong di bawah ambangnya dan kromatinnya dibuka kembali sekaligus,
yang hanya dapat dikerjakan sitoplasma telur atau empat faktor pada
sebagian kecil selnya: mungkin, karena tidak ada yang hilang, tetapi
tidak berdaya guna, karena dua gembok yang saling bebas harus dicungkil
bersama-sama.
\end{solution}

\begin{solution}{pb:b2:cell-differentiation:1}
\textbf{1.} $0.2/2.5\times 10^{-6} = \num{80000}$ \omterm{def:b2:cell-differentiation:fibre}{sarkomer};
$0.2/25\times 10^{-6} = \num{8000}$ inti.
\textbf{2.} Kira-kira \num{8000} tiap serabutnya; $4\times 10^{9}$ bagi
ototnya.
\textbf{3.} $\pi(30\times 10^{-6})^{2}\times 0.2 = \qty{5.7e-10}{m^3}$
tiap serabutnya; $2.8\times 10^{-4}$ \unit{m^3}, yaitu \qty{0.28}{L},
bagi ototnya.
\textbf{4.} $5.7\times 10^{-10}/8000 = \qty{7e-14}{m^3} =
\qty{70000}{\micro m^3}$: yaitu setara 35 sel biasa tiap intinya.
\textbf{5.} Bagian dalamnya terjejali \omterm{def:b2:cell-differentiation:fibre}{miofibril} yang harus berjalan
tanpa terputus dari ujung ke ujung; inti di tepinya membiarkan kisinya
tetap sinambung dan duduk dekat membrannya beserta isyaratnya.
\textbf{6.} $4\times 10^{9}/2000 = 2\times 10^{6}$: yaitu 21 kali
pelipatduaan, 10,5 hari.
\textbf{7.} $0.4m = m^{2}/(1 + m^{2}) + 0.02$. Pada $m = 0.055$: kiri
$0.022$, kanan $0.003 + 0.02 = 0.023$.
\textbf{8.} Pada $m = 2.1$: kiri $0.84$, kanan $0.815 + 0.02 = 0.835$.
\textbf{9.} Pada $m = 0.41$: kiri $0.164$, kanan $0.144 + 0.02 = 0.164$.
Tepat di bawahnya sintesisnya lebih kecil daripada penguraiannya,
sehingga $m$ turun menuju $0.055$; tepat di atasnya, sintesisnya
melampaui penguraiannya dan $m$ naik menuju $2.1$: tiap simpangannya
membesar --- tak mantap.
\textbf{10.} $0.6 > 0.41$: sel pertamanya pergi ke keadaan nyala lalu
berdiferensiasi; $0.3 < 0.41$: sel keduanya mengendur ke $0.055$ dan
tidak berdiferensiasi.
\textbf{11.} Dengan $b = 0.4$ kurva sintesisnya bermula pada $0.4$ lalu
tinggal di atas garis $0.4m$ sampai $m \approx 3.3$: perpotongan
rendahnya lenyap. Selnya berdiferensiasi secara spontan, tanpa
pengimbasan.
\textbf{12.} Pembelahannya memerohkan $m$; anaknya mewarisi $2.1/2 =
1.05$, yang di atas ambangnya $0.41$, lalu memanjat kembali ke $2.1$:
keadaannya dibangun kembali di dalam tiap anaknya. Pemancangannya
selamat selama kadar yang terpeorh itu tinggal di atas ambangnya.
\textbf{13.} $4$ tiap mm $\times$ \qty{200}{mm} $= 800$ sel satelit
tiap serabutnya; \qty{10}{\%} dari \num{8000} intinya, yaitu \num{800},
harus diganti.
\textbf{14.} Ruas yang cedera memuat $80$ sel satelit; $800/80 =
10$,
sehingga 4 kali pelipatduaan ($\times 16$, yang memberi \num{1280}: 800
melebur, 480 tersisa untuk memulihkan lumbungnya) --- yaitu
\qty{72}{h}, tiga hari.
\textbf{15.} Tiga hari pembelahan terhadap dua sampai tiga pekan yang
teramati: sisanya adalah pembersihan puingnya oleh makrofag, jeda
pengaktifannya, peleburannya, perakitan \omterm{def:b2:cell-differentiation:fibre}{miofibrilnya}, penyarafan
ulangnya, dan pematangan ruas barunya.
\textbf{16.} \omterm{def:b2:cell-differentiation:myogenesis}{MyoD} (dan miogenin) harus menyalakan kembali gen ototnya
di dalam sel yang diaktifkan; pembelahan dan \omterm{def:b2:cell-differentiation:differentiation}{diferensiasi} saling
meniadakan --- Id-nya harus turun dan daurnya harus berhenti sebelum
miogeninnya dapat bekerja dan peleburannya terjadi.
\textbf{17.} $0.95^{n} = 0.5$: $n = 13.5$, yaitu kira-kira empat belas
cedera.
\textbf{18.} $800\times 0.95^{20} = 800\times 0.36 \approx 290$.
\textbf{19.} Tiap pengerutannya merobek serabutnya, tiap robekannya
memanggil lumbungnya, dan lumbungnya mengerut beberapa persen tiap
kalinya; sesudah bertahun-tahun ia habis, serabut yang robek tidak
dibangun kembali, dan fibroblas mengisi celahnya dengan kolagen.
\textbf{20.} $(80/60)^{2} = 1.78$; intinya $8000\times 1.78 = \num{14200}$:
yaitu kira-kira \num{6200} tambahan dari sel satelitnya.
\textbf{21.} $0.28\times 1.78 = \qty{0.5}{L}$.
\textbf{22.} Yang berubah: gen isoform miosinnya, kerapatan mitokondria
dan kapilernya, mioglobinnya. Yang tidak berubah: jati diri ototnya
(keadaan \omterm{def:b2:cell-differentiation:myogenesis}{MyoD-nya}), arsitektur \omterm{def:b2:cell-differentiation:fibre}{sarkomernya}, intinya, dan genomnya.
\textbf{23.} Serabutnya pascamitosis dan tidak dapat berlipat ganda;
satu-satunya jalannya adalah serabut yang lebih tebal dan inti yang
lebih banyak tiap serabutnya dari sel satelitnya.
\textbf{24.} Latihannya: penebalan yang terbatas, karena serabutnya
tidak dapat menambah inti; cederanya: tanpa perbaikan --- ruas yang
musnah digantikan jaringan berserat, seperti pada distrofi.
\textbf{25.} \num{8000} inti tiap serabutnya; ambangnya $m \approx 0.41$;
kira-kira tiga hari pembelahan sel satelit bagi sebuah perbaikan.
\end{solution}
